\section[Cote 133 : le centre et le groupe dérivé vus par deux modules croisés]{Cote 133 : le centre et le groupe dérivé vus par deux modules croisés\protect\verifie{Folder133}}\label{sec:133}

La cote 133, \emph{[Groupes de Witt et formes quadratiques, groupes formels] : notes manuscrites (s.d.), tirés à part (1974)}, que l'inventaire date \q{1974-[à partir de 1975]}, réunit sept suites distinctes. La dernière, sous une chemise titrée \q{Fourbis auxiliaires cohomologiques pour groupes formels} (page~44), occupe les pages~45 à~53, qu'il pagine de 1 à 7 ; les pages~47 et~49 ne sont pas de sa main. Elle traite en général le problème que la \q{Solution} des pages~35 et~36 pose pour un groupe formel : décrire un groupe $G$ par son centre et son groupe dérivé. La page~45 se donne deux modules croisés $\mathcal C = (N \to M)$ et $\mathcal C' = (N' \to M')$ de mêmes invariants $\pi_0$, $\pi_1$, et note $K = G/\pi_1$ (\lot{133}{3}{45}). La page~46 définit $\mathfrak Z = [\mathrm{Cent}(K)/(N'/\pi_1)] \cap \operatorname{Ker}\Theta$, construit $\varphi_G$ puis $\Psi_G$, et encadre \q{$\Psi_G$ injectif} comme équivalent à $N' = \mathrm{Cent}(G)$ (\lot{133}{3}{46}). La page~48 définit l'accouplement $c_E$ d'une extension centrale $E$ de $\pi_0$ par $\pi_1$ et écrit \q{On doit trouver, au signe près} $\Psi_{G_1} = \Psi_G + \tilde c_E\,\alpha$ (\lot{133}{3}{48}). Les pages~51 et~52 définissent $\mathcal D = (N_{\mathrm{ab}})_M$ et $\lambda_G$, et écrivent $N/DG \simeq \mathcal D/\lambda_G(\pi_0 \otimes \pi_0)$ et $\lambda_{G_1} = \lambda_G + \beta c_E$ (\lot{133}{3}{51}, \lot{133}{3}{52}). Malgré le titre, rien n'y est propre aux groupes formels, et les énoncés sont démontrés ici pour des groupes abstraits.

\noindent\emph{Conventions.} Le commutateur est $[x, y] = xyx^{-1}y^{-1}$. Dans tout ce qui suit, $G$ est un groupe, $N$ et $N'$ deux sous-groupes distingués, avec $N' \subset \mathrm{Cent}(G)$ ; alors $[N, N'] = 1$, $\pi_1 = N \cap N'$ est central, $\pi_0 = G/NN'$, $K = G/\pi_1$, $M = G/N'$ agit sur $N$ par conjugaison, et $\Theta : M \to \mathrm{Aut}(N)$ est cette action. Le groupe $\mathfrak Z$ est pris par son image inverse $\overline{N'}$ dans $G$ : les $z \in G$ dont l'image dans $K$ est centrale et qui commutent à tout élément de $N$, c'est-à-dire tels que $[g, z] \in \pi_1$ pour tout $g \in G$ et $nz = zn$ pour tout $n \in N$. On a $N' \subset \overline{N'}$ et $\mathfrak Z = \overline{N'}/N'$.

\begin{lemme}\label{lem:133-central}
Si $c$ est central dans un groupe, $[g, zc] = [g, z]$ et $[gc, z] = [g, z]$.
\end{lemme}

\begin{proof}
$[g, zc] = gzcg^{-1}c^{-1}z^{-1} = gzg^{-1}cc^{-1}z^{-1} = [g, z]$, en faisant passer $c$ devant $g^{-1}$. La seconde s'en déduit par $[x, y]^{-1} = [y, x]$.
\end{proof}

\begin{enonce}[Page 46, l'injectivité de $\Psi_G$\piste{133-centre-and-derived-group-from-two-crossed-modules}]
Pour $z \in \overline{N'}$, l'application $g \mapsto [g, z]$ définit un homomorphisme $\varphi_G(z) : \pi_0 \to \pi_1$, qui ne dépend que de la classe de $z$ dans $\mathfrak Z$ ; l'application $\Psi_G : \mathfrak Z \to \mathrm{Hom}(\pi_0, \pi_1)$ ainsi définie est un homomorphisme, et
\[
N' = \mathrm{Cent}(G) \iff \Psi_G \text{ est injectif.}
\]
\end{enonce}

\begin{proof}
Soit $z \in \overline{N'}$. Par définition $[g, z] \in \pi_1$, donc $[g, z]$ est central. On a $[gh, z] = g\,[h, z]\,zg^{-1}z^{-1}$ ; comme $[h, z]$ est central, c'est $[h, z]\,[g, z] = [g, z]\,[h, z]$, et $g \mapsto [g, z]$ est un homomorphisme $G \to \pi_1$. Il est trivial sur $N$, puisque $z$ commute à $N$, et sur $N'$, qui est central ; il se factorise donc par $\pi_0 = G/NN'$, ce qui définit $\varphi_G(z)$. Si $z$ et $w$ ont même classe modulo $N'$, $w = zc$ avec $c \in N'$ central, et $[g, w] = [g, z]$ par le lemme~\ref{lem:133-central} : d'où $\Psi_G$ sur $\mathfrak Z$. Il est multiplicatif : $[g, zw] = gzg^{-1}\,[g, w]\,z^{-1} = [g, z]\,[g, w]$, puisque $[g, w]$ est central.

Supposons $\Psi_G$ injectif et soit $c \in \mathrm{Cent}(G)$. Alors $c \in \overline{N'}$, et $\Psi_G(\bar c)$ est trivial, comme $\Psi_G(1)$ : donc $\bar c = 1$ dans $\mathfrak Z$, c'est-à-dire $c \in N'$. Ainsi $\mathrm{Cent}(G) \subset N'$, et l'égalité vient de $N' \subset \mathrm{Cent}(G)$.

Réciproquement, supposons $N' = \mathrm{Cent}(G)$, et soient $z, w \in \overline{N'}$ avec $\Psi_G(\bar z) = \Psi_G(\bar w)$, c'est-à-dire $[g, z] = [g, w]$ pour tout $g$. Appliquée à $g^{-1}$, l'égalité donne $zgz^{-1} = wgw^{-1}$, d'où $g\,z^{-1}w = z^{-1}(zgz^{-1})w = z^{-1}(wgw^{-1})w = z^{-1}w\,g$ : $z^{-1}w$ est central, donc dans $N'$, et $\bar z = \bar w$.
\end{proof}

\begin{remarque}
Seule l'hypothèse $N' \subset \mathrm{Cent}(G)$ sert. L'hypothèse $[N, N'] = 1$ que la lecture ajoute en découle, et $N \supset D(G)$ n'intervient pas ici. La vérification confirme les contrôles d'additivité et de nullité que la lecture fournit et que la page résume par \q{triviale sur $N'$ et sur $\mathfrak Z$}. Elle établit en chemin que $\mathrm{Cent}(G)$ est l'image inverse de $\operatorname{Ker}\Psi_G$, ce que la page dit \q{clair}.
\end{remarque}

On suppose maintenant de plus $D(G) \subset N$. Soit $\mathcal D = (N_{\mathrm{ab}})_M$ ; c'est le quotient $N/[G, N]$, le sous-groupe $[G, N]$ engendré par les $[g, n]$ contenant les $[n, n']$. On note $\bar n$ la classe dans $\mathcal D$ d'un $n \in N$.

\begin{lemme}\label{lem:133-D}
Le groupe $\mathcal D$ est commutatif, $G$ y agit trivialement, et $\pi_0$ est commutatif.
\end{lemme}

\begin{proof}
Pour $a, b \in N$, $(ab)^{-1}(ba) = [b^{-1}, a^{-1}] \in [G, N]$, donc $\bar a\bar b = \bar b\bar a$. Pour $c \in N$ et $x \in G$, $(xcx^{-1})^{-1}c = [x, c^{-1}] \in [G, N]$, donc $\overline{xcx^{-1}} = \bar c$. Pour $x, y \in G$, $(xy)^{-1}(yx) = [y^{-1}, x^{-1}] \in D(G) \subset N \subset NN'$, donc $xy$ et $yx$ ont même image dans $\pi_0$.
\end{proof}

\begin{enonce}[Pages 51 et 52, la surjectivité de $\lambda_G$\piste{133-centre-and-derived-group-from-two-crossed-modules}]
Supposons $D(G) \subset N$ et $N' \subset \mathrm{Cent}(G)$. L'application $(x, y) \mapsto \overline{[x, y]}$ définit une forme $\lambda_G : \pi_0 \times \pi_0 \to \mathcal D$, multiplicative en chaque variable et alternée, et
\[
N = D(G) \iff \lambda_G \text{ est surjectif,}
\]
c'est-à-dire si et seulement si les valeurs de $\lambda_G$ engendrent $\mathcal D$.
\end{enonce}

\begin{proof}
Pour $y$ fixé, $x \mapsto \overline{[x, y]}$ est un homomorphisme $G \to \mathcal D$ : $[xx', y] = (x[x', y]x^{-1})[x, y]$, et le lemme~\ref{lem:133-D} donne $\overline{[xx', y]} = \overline{[x', y]}\;\overline{[x, y]} = \overline{[x, y]}\;\overline{[x', y]}$. Il est trivial sur $N$, car $[n, y] = [y, n]^{-1} \in [G, N]$, et sur $N'$, car $[n', y] = 1$ pour $n'$ central ; il se factorise donc par $\pi_0$. Comme $\overline{[x, y]} = \overline{[y, x]}^{-1}$, la même chose vaut pour la seconde variable, et $\lambda_G$ est bien défini sur $\pi_0 \times \pi_0$, multiplicatif en chaque variable ; $[x, x] = 1$ le rend alterné.

Si $N = D(G)$, tout $n \in N$ est produit de commutateurs et de leurs inverses ; sa classe $\bar n$ est le produit correspondant de valeurs de $\lambda_G$ et de leurs inverses, et ces valeurs engendrent $\mathcal D$. Réciproquement, les valeurs de $\lambda_G$ sont des classes d'éléments de $D(G)$, donc le sous-groupe qu'elles engendrent est contenu dans l'image de $D(G)$ dans $\mathcal D$. S'il est égal à $\mathcal D$, tout $n \in N$ s'écrit $n = m\,(m^{-1}n)$ avec $m \in D(G)$ et $m^{-1}n \in [G, N] \subset D(G)$ ; donc $N \subset D(G)$, et $N = D(G)$.
\end{proof}

\begin{remarque}
Comme $\mathcal D$ est commutatif et $\lambda_G$ bimultiplicatif, $\lambda_G$ définit une application linéaire $\pi_0 \otimes \pi_0 \to \mathcal D$, dont l'image est le sous-groupe engendré par les valeurs de $\lambda_G$ ; c'est la surjectivité de la page. L'hypothèse $N' \subset \mathrm{Cent}(G)$ est nécessaire à la définition même de $\lambda_G$ sur $\pi_0$, et la lecture la cite (\q{puisque $N'$ est central}). L'hypothèse plus faible $[N, N'] = 1$ ne suffit pas à l'équivalence. Dans le groupe des quaternions $Q_8$, avec $N = D(Q_8) = \mathrm{Cent}(Q_8)$, d'ordre~$2$, et $N' = Q_8$, on a $[N, N'] = 1$ et $\pi_0 = 1$, donc toute forme sur $\pi_0$ est nulle, alors que $\mathcal D = N/[G, N] = N$ n'est pas trivial (exemple non vérifié en Lean). La commutativité de $\pi_0$ découle de $D(G) \subset N$ (lemme~\ref{lem:133-D}) et n'est pas une hypothèse de plus. L'isomorphisme $N/D(G) \simeq \mathcal D/\lambda_G(\pi_0 \otimes \pi_0)$ de la page n'est démontré ici que sous la forme de l'équivalence.
\end{remarque}

Soient enfin $E$ un groupe, $\iota : \pi_1 \to E$ un homomorphisme injectif d'image centrale, et $q : E \to \pi_0$ un homomorphisme de noyau $\iota(\pi_1)$ : une extension centrale de $\pi_0$ par $\pi_1$, si $q$ est surjectif, ce dont on ne se sert pas. Soit $G \times_{\pi_0} E$ le produit fibré et $G_1$ son quotient par l'image antidiagonale $\{(a, \iota(a)^{-1})\}$ de $\pi_1$, centrale, prise en Lean par son sous-groupe distingué engendré.

\begin{enonce}[Pages 48 et 52, le changement de $G$ en $G_1$\piste{133-centre-and-derived-group-from-two-crossed-modules}]
Supposons $D(G) \subset N$ et $N' \subset \mathrm{Cent}(G)$.
\begin{enumerate}
\item Pour $e, f \in E$, $[e, f] = \iota(c_E(e, f))$ pour un unique $c_E(e, f) \in \pi_1$, qui ne dépend que des images de $e$ et $f$ dans $\pi_0$.
\item Dans $G_1$, le commutateur des classes de $(g, e)$ et $(z, f)$ est la classe de $([g, z]\,c_E(e, f), 1)$.
\item En particulier, si $z \in \overline{N'}$, ce commutateur est l'image de $\Psi_G(\bar z)(\bar g)\,c_E(e, f) \in \pi_1$ ; et la classe dans $\mathcal D$ de $[x, y]\,c$, pour $c \in \pi_1$, est $\lambda_G(\bar x, \bar y)\,\beta(c)$, où $\beta : \pi_1 \to \mathcal D$ est le passage au quotient.
\end{enumerate}
\end{enonce}

\begin{proof}
(1) L'image de $[e, f]$ dans $\pi_0$ est le commutateur des images, trivial puisque $\pi_0$ est commutatif ; $[e, f]$ est donc dans $\operatorname{Ker} q = \iota(\pi_1)$, et $\iota$ est injectif. Si $e$ et $e'$ ont même image, $e' = e\,\iota(a)$, et de même $f' = f\,\iota(b)$ ; comme $\iota(a)$ et $\iota(b)$ sont centraux, le lemme~\ref{lem:133-central} donne $[e', f'] = [e, f]$.

(2) Le commutateur de $(g, e)$ et $(z, f)$ dans le produit fibré est $([g, z], [e, f]) = ([g, z], \iota(c))$, $c = c_E(e, f)$, et
\[
([g, z], \iota(c))^{-1}\,([g, z]\,c, 1) = (c, \iota(c)^{-1})
\]
est dans l'image antidiagonale de $\pi_1$ ; les deux éléments ont donc même classe dans $G_1$. L'élément $([g, z]\,c, 1)$ est bien dans le produit fibré, $[g, z]\,c$ étant dans $N$.

(3) Si $z \in \overline{N'}$, $[g, z] = \Psi_G(\bar z)(\bar g)$ par définition. La seconde assertion est $\overline{[x, y]\,c} = \overline{[x, y]}\;\bar c$.
\end{proof}

\begin{remarque}
Ce calcul porte les deux formules de la page, $\Psi_{G_1} = \Psi_G + \tilde c_E\,\alpha$ et $\lambda_{G_1} = \lambda_G + \beta c_E$, écrites additivement : dans $G_1$, un commutateur est le produit du commutateur dans $G$ et de la valeur de $c_E$. L'ordre des facteurs fixe le signe que la page laisse ouvert (\q{au signe près}). Avec $[x, y] = xyx^{-1}y^{-1}$ et $\Psi_G(\bar z)(\bar g) = [g, z]$, le terme correctif est $c_E(\bar g, \bar z)$, c'est-à-dire $\tilde c_E(\alpha(\bar z))$ évalué en $\bar g$, au signe près selon le côté où l'on curryfie. L'identification des deux modules croisés de $G_1$ avec ceux de $G$ n'est pas formalisée ; elle ferait de cette égalité d'éléments une égalité d'applications $\Psi_{G_1} = \Psi_G \cdot \tilde c_E\,\alpha$ sur un même $\mathfrak Z$.
\end{remarque}
